% !TeX root = ../main.tex
% ============================================================
%  sections/03_problem_description.tex
% ============================================================

\section{Problem Description}
\label{sec:problem}

We study the tactical design of a multi-echelon, multi-period \CLSC
operating under an external carbon-tax policy. The resulting decision
problem is hierarchical: a regulator sets the tax rate, and a supply
chain operator designs and operates the network under that rate. \Cref{sec:prob:overview}
formalizes the physical network and admissible flows;
\Cref{sec:prob:hierarchy} states the Stackelberg interaction that
underpins the bilevel model of \Cref{sec:model}.

%-----------------------------------------------------------------
\subsection{Network and flows}
\label{sec:prob:overview}

The planning horizon comprises $T$ discrete periods. The physical network
consists of five node classes: suppliers ($\setS$), factories ($\setF$),
warehouses ($\setW$), hub terminals $\setH=\setHa\cup\setHs$ (airports and
seaports), and markets ($\setR$). Suppliers, factories, warehouses, and
hub terminals are \emph{openable}. Opening a factory or warehouse incurs
a construction cost and an emission, both proportional to its
installed area -- there is no fixed cost or emission for opening a
factory or warehouse beyond what its installed area requires.
Suppliers and hub terminals are not constructed but
\emph{selected}: selecting a supplier incurs a fixed cost, reflecting
the footprint of qualifying a new sourcing relationship rather than
building a facility, but no emission of its own -- a supplier's
emission footprint is captured entirely by the raw-material
acquisition emission charged on what is actually shipped from it --
whereas selecting a hub incurs a fixed access cost but no construction
emission, since hub terminals are modeled as outsourced terminal
access points. Markets face
demand given as a model parameter and act as collection points for
end-of-use items.

Three item families circulate in the network. Raw materials ($\setU$)
are procured from suppliers and consumed at factories. Finished products
($\setP$) are manufactured or remanufactured at factories and delivered
to markets through warehouses, hubs, or direct links. Recoverable items
($\setK$) are collected at markets and returned upstream for
remanufacturing, closing the material loop. Forward and reverse flows
share the same transportation infrastructure.

Transportation combines a heterogeneous modal portfolio. Road movement
uses company-owned truck alternatives, so the model can choose among
truck types with different capacities, fixed costs, trip costs, and
emission intensities. Long-haul inter-hub transfers use outsourced
airplane and boat alternatives, represented by carrier contracts with per-period
capacity, variable kg--km costs, and terminal handling costs. Hub
terminals operate as intermodal cross-docking facilities: ground inflows
are matched by air or sea outflows of the same tonnage, and vice versa.

Production and remanufacturing at each factory rely on a catalogue of
technologies ($\setG$ and $\setL$, respectively), each characterized by
unit cost, throughput bounds, emission factor, and a set of compatible
products. Finished-product inventories are held at factories and
warehouses subject to opening-conditional capacity limits; markets
satisfy demand exactly with no carry-over. Returned items enter the
reverse channel with a product-specific lag and must be absorbed by
remanufacturing at a factory.

\Cref{fig:network_structure} depicts the admissible arc families.
Solid arcs denote forward flows (direct and warehouse-mediated), dashed
arcs denote reverse flows, and grey arcs denote intermodal routes
through the hub layer. The two hub nodes in the figure represent the
same set $\setH$ in its origin and destination roles along an
intermodal path. Demand and return-rate bounds are deterministic and
known at planning time, consistent with the tactical scope of the
problem.

\begin{figure}[htbp]
  \centering
  \begin{tikzpicture}[
      >=Latex,
      node distance=2.8cm,
      every node/.style={align=center, font=\small},
      % --- node styles: grayscale-friendly, tonal differentiation via fill ---
      facility/.style={
        draw=black!75, line width=0.7pt, fill=black!4,
        rectangle, rounded corners=3pt,
        minimum width=1.45cm, minimum height=0.85cm
      },
      hub/.style={
        draw=black!75, line width=0.7pt, fill=black!12,
        ellipse, minimum width=1.8cm, minimum height=0.95cm
      },
      % --- arc styles: three line types, one colour, print-safe ---
      fwd/.style ={-{Latex[length=2.0mm]}, line width=0.9pt, draw=black!85},
      rev/.style ={-{Latex[length=2.0mm]}, line width=0.9pt, draw=black!85,
                   dashed},
      imod/.style={-{Latex[length=1.6mm]}, line width=0.55pt, draw=black!45,
                   densely dotted},
      % --- label styles ---
      flab/.style={font=\small, fill=white, inner sep=1.3pt},
      role/.style={font=\footnotesize\itshape, text=black!55}
    ]
    % ================= Main chain =================
    \node[facility]             (S) {$\setS$};
    \node[facility, right=of S] (F) {$\setF$};
    \node[facility, right=of F] (W) {$\setW$};
    \node[facility, right=of W] (R) {$\setR$};
    % ================= Hub layer (two routing roles) ==
    \node[hub, above=2.7cm of F, xshift=-0.35cm] (Ho) {$\setH$};
    \node[hub, above=2.7cm of W, xshift=0.35cm]  (Hd) {$\setH$};
    \node[role, above=2pt of Ho] {origin role};
    \node[role, above=2pt of Hd] {destination role};
    % ================= Forward flows (solid, above chain) =====
    \draw[fwd] (S) -- node[flab, above] {raw materials} (F);
    \draw[fwd] (F) -- node[flab, above] {products}      (W);
    \draw[fwd] (W) -- node[flab, above] {delivery}      (R);
    \draw[fwd] (F) to[bend left=14]
               node[flab, pos=0.42, above=2pt] {direct delivery} (R);
    % ================= Reverse flows (dashed, below chain) =====
    \draw[rev] (W) to[bend left=16]
               node[flab, pos=0.38, below=1pt] {recovered items}  (F);
    \draw[rev] (R) to[bend left=16]
               node[flab, pos=0.62, below=1pt] {returns}          (W);
    \draw[rev] (R) to[bend left=38]
               node[flab, pos=0.5, below=3pt] {direct returns} (F);
    % ================= Intermodal: into origin-role hub (S, F, W, R) =====
    \draw[imod] (S.north) to[out=85,  in=210] (Ho.south west);
    \draw[imod] (F.north) to[out=95,  in=255] (Ho.south);
    \draw[imod] (W.north) .. controls ($(W)+(-0.25,1.6)$)
                          and ($(Ho)+(0.95,-0.75)$) .. (Ho.east);
    \draw[imod] (R.north) .. controls ($(R)+(-0.5,1.9)$)
                          and ($(Ho)+(0.85,-0.9)$) .. (Ho.south east);
    % ================= Intermodal: out of destination-role hub (F, W, R) =
    \draw[imod] (Hd.south west) .. controls ($(Hd)+(-0.85,-0.9)$)
                                and ($(F)+(0.15,1.9)$) .. (F.north);
    \draw[imod] (Hd.south) -- (W.north);
    \draw[imod] (Hd.south east) to[out=290, in=85] (R.north);
    % ================= Hub transfer =============================
    \draw[imod] (Ho) -- node[flab, above]     {air / sea transfer}
                        node[role, below=5pt] {(same set $\setH$)} (Hd);
    % ================= Legend =========================
    \begin{scope}[shift={([yshift=-2.2cm]S.south west)}]
      \draw[fwd]  (0,0)   -- ++(0.9,0)
        node[right, font=\small, inner sep=3pt] {forward flow};
      \draw[rev]  (3.0,0) -- ++(0.9,0)
        node[right, font=\small, inner sep=3pt] {reverse flow};
      \draw[imod] (6.0,0) -- ++(0.9,0)
        node[right, font=\small, inner sep=3pt] {intermodal routing};
    \end{scope}
  \end{tikzpicture}
  \caption{Forward (solid), reverse (dashed), and intermodal (dotted)
    flows in the \CLSC network, connecting suppliers ($\setS$),
    factories ($\setF$), warehouses ($\setW$), and markets ($\setR$).
    Intermodal routes pass through the hub layer $\setH$, shown twice
    to distinguish its origin and destination roles.}
  \label{fig:network_structure}
\end{figure}

%-----------------------------------------------------------------
\subsection{Bilevel structure}
\label{sec:prob:hierarchy}

Network design and operation are governed by a Stackelberg game between
a regulator and a supply chain operator: the regulator commits to a
carbon tax before the operator observes it and responds.

\subsubsection{Leader (regulator)}
\label{sec:prob:leader}

The regulator chooses the carbon tax
from a finite, discrete grid of rates, which enables the exact
decomposition of \Cref{sec:methodology}. The regulator's objective is
\emph{lexicographic}: first minimize the supply chain's total
emissions, and, among all tax rates that attain that minimum, levy
the lowest one.

\subsubsection{Follower (supply chain operator)}
\label{sec:prob:follower}

Given the announced tax
rate, the operator minimizes total operating cost plus tax payments
by jointly deciding
facility opening and installed capacity, technology installation,
transportation modes and fleet, forward and reverse product flows,
production, remanufacturing, and inventory.
